% This file was converted to LaTeX by Writer2LaTeX ver. 1.9.9
% see http://writer2latex.sourceforge.net for more info
\documentclass{article}
\usepackage{calc,amsmath,amssymb,amsfonts}
\usepackage[T1]{fontenc}
\usepackage[italian]{babel}
\usepackage[style=numeric,backend=biber]{biblatex}
\usepackage{tikz}
\author{HP}
\date{2026-07-26}
\begin{document}
\section{Prompt editoriale – Progetto del libro}
\subsection{Titolo provvisorio}
Come sviluppare un videogame con il Commodore 64 all'epoca di ChatGPT

Sottotitolo da definire durante il progetto.

\begin{tikzpicture}
\path  (0,0.054) --  (17.001,0.054) --  (17.001,0) --  (0,0) --  (0,0.054) -- cycle;
\end{tikzpicture}


\section{Visione del progetto}
Questo libro nasce da un'idea semplice ma ambiziosa: riportare il lettore davanti al Commodore 64 con lo stesso
entusiasmo che provava negli anni Ottanta, accompagnandolo però con gli strumenti del XXI secolo.

Non deve essere un semplice manuale tecnico.

Non deve essere un libro nostalgico.

Non deve essere un corso di BASIC.

Non deve essere un corso di Assembly.

Deve essere un viaggio.

Il lettore non deve limitarsi a imparare come funziona il Commodore 64.

Deve tornare a sentirsi il ragazzo curioso che passava interi pomeriggi a sperimentare davanti a uno schermo blu.

\begin{tikzpicture}
\path  (0,0.054) --  (17.001,0.054) --  (17.001,0) --  (0,0) --  (0,0.054) -- cycle;
\end{tikzpicture}


\section{Il lettore}
Il libro è pensato contemporaneamente per tre categorie di persone.

Il primo è l'ex ragazzo degli anni Ottanta che ritrova il proprio Commodore 64 in cantina e decide di riaccenderlo dopo
quarant'anni.

Il secondo è il programmatore moderno che preferisce utilizzare un emulatore e desidera comprendere come venivano
realizzati i videogiochi dell'epoca.

Il terzo è il giovane appassionato che non ha mai posseduto un Commodore 64 ma vuole capire perché quella macchina abbia
cambiato la storia dell'informatica personale.

Il libro deve riuscire a parlare a tutti e tre senza privilegiare nessuno.

\begin{tikzpicture}
\path  (0,0.054) --  (17.001,0.054) --  (17.001,0) --  (0,0) --  (0,0.054) -- cycle;
\end{tikzpicture}


\section{Filosofia editoriale}
Il videogioco rappresenta il filo conduttore dell'intero libro.

Non è il fine.

È il mezzo attraverso il quale vengono introdotti tutti gli argomenti.

Ogni capitolo deve soddisfare contemporaneamente quattro obiettivi.

Far avanzare la storia.

Far crescere concretamente il videogioco.

Introdurre almeno un nuovo concetto tecnico.

Far rivivere al lettore un'emozione o una scoperta tipica degli anni Ottanta.

La teoria non deve mai precedere il bisogno.

Ogni spiegazione deve nascere da un problema reale incontrato durante lo sviluppo del videogioco.

Il lettore non deve mai avere la sensazione di studiare.

Deve avere l'impressione di costruire qualcosa.

\begin{tikzpicture}
\path  (0,0.054) --  (17.001,0.054) --  (17.001,0) --  (0,0) --  (0,0.054) -- cycle;
\end{tikzpicture}


\section{Filosofia narrativa}
Lo stile deve essere quello di un romanzo tecnico.

La narrazione deve essere discorsiva, coinvolgente e cinematografica.

Il racconto deve fondersi continuamente con la spiegazione tecnica.

Le emozioni non devono interrompere la teoria.

La teoria non deve interrompere la narrazione.

L'intero libro deve essere scritto come una lunga storia.

Ogni capitolo deve concludersi lasciando il desiderio di leggere quello successivo.

\begin{tikzpicture}
\path  (0,0.054) --  (17.001,0.054) --  (17.001,0) --  (0,0) --  (0,0.054) -- cycle;
\end{tikzpicture}


\section{Il ruolo di ChatGPT}
ChatGPT è uno dei due coautori del libro.

La sua presenza deve essere dichiarata fin dalla prefazione.

Non rappresenta il protagonista.

Non sostituisce il programmatore.

Non scrive il videogioco al posto del lettore.

Assume il ruolo del compagno di laboratorio.

È l'equivalente moderno dell'amico più esperto che, negli anni Ottanta, suggeriva una soluzione, consigliava una rivista
o spiegava un comando misterioso.

L'intelligenza artificiale diventa uno strumento di apprendimento, non un sostituto della creatività.

\begin{tikzpicture}
\path  (0,0.054) --  (17.001,0.054) --  (17.001,0) --  (0,0) --  (0,0.054) -- cycle;
\end{tikzpicture}


\section{Il protagonista}
Il vero protagonista del libro è il lettore.

Il Commodore 64 rappresenta il laboratorio.

Il videogioco rappresenta il progetto.

ChatGPT rappresenta il compagno di viaggio.

\begin{tikzpicture}
\path  (0,0.054) --  (17.001,0.054) --  (17.001,0) --  (0,0) --  (0,0.054) -- cycle;
\end{tikzpicture}


\section{Il videogioco}
Il videogioco costituisce l'intera struttura del libro.

Ogni capitolo aggiunge realmente nuove funzionalità.

Il progetto cresce in modo progressivo.

Il lettore deve poter eseguire il codice alla fine di ogni capitolo.

Il protagonista del gioco sarà un drone.

Il drone costituisce una reinterpretazione moderna della celebre mongolfiera presente nel manuale originale del
Commodore 64.

Il drone rappresenta il collegamento simbolico tra passato e futuro.

Il genere del videogioco dovrà essere scelto privilegiando la gradualità didattica, permettendo l'introduzione naturale
di sprite, scrolling, raster interrupt, collisioni, suoni, musica, intelligenza artificiale elementare e ottimizzazioni
in Assembly.

\begin{tikzpicture}
\path  (0,0.054) --  (17.001,0.054) --  (17.001,0) --  (0,0) --  (0,0.054) -- cycle;
\end{tikzpicture}


\section{Progressione tecnica}
Il libro insegna contemporaneamente:

\begin{itemize}
\item BASIC del Commodore 64; 
\item Assembly del MOS 6510; 
\item architettura del Commodore 64; 
\item memoria; 
\item VIC-II; 
\item SID; 
\item sprite; 
\item raster interrupt; 
\item scrolling; 
\item collisioni; 
\item tecniche di game development; 
\item tecniche di ottimizzazione; 
\item metodo di progettazione dei programmatori degli anni Ottanta. 
\end{itemize}
Gli argomenti non devono seguire un ordine accademico.

Devono seguire esclusivamente le necessità del videogioco.

L'Assembly compare soltanto quando il lettore comprende naturalmente i limiti del BASIC.

\begin{tikzpicture}
\path  (0,0.054) --  (17.001,0.054) --  (17.001,0) --  (0,0) --  (0,0.054) -- cycle;
\end{tikzpicture}


\section{Prima di iniziare il viaggio}
Prima dell'inizio del libro è prevista una sezione introduttiva dedicata alla preparazione dell'esperienza.

Questa sezione accompagnerà il lettore nella scelta del percorso più adatto.

Per chi possiede ancora un Commodore 64, verranno illustrate le precauzioni necessarie prima della riaccensione dopo
molti anni: controllo dell'alimentatore, verifiche di sicurezza, pulizia, consigli per evitare danni e indicazioni su
quando sia preferibile rivolgersi a un tecnico specializzato.

Per chi preferisce lavorare in emulazione, verrà spiegata l'installazione e la configurazione di VICE, la preparazione
dell'ambiente di sviluppo moderno e l'organizzazione degli strumenti utilizzati durante tutto il libro.

Entrambi i percorsi dovranno convergere rapidamente, affinché nessun lettore si senta escluso.

Durante il resto del libro sarà chiaramente indicato quando un'esperienza è specifica del Commodore reale,
dell'emulatore o comune a entrambi.

\begin{tikzpicture}
\path  (0,0.054) --  (17.001,0.054) --  (17.001,0) --  (0,0) --  (0,0.054) -- cycle;
\end{tikzpicture}


\section{Materiale didattico}
Ogni capitolo dovrà contenere, quando opportuno, rubriche ricorrenti.

Laboratorio

Realizzazione pratica di una nuova parte del videogioco.

Alla fine di ogni laboratorio il progetto cresce realmente.

Esperimento

Piccole attività guidate in cui il lettore modifica il codice, osserva gli effetti e comprende i concetti attraverso la
sperimentazione.

L'errore è parte integrante dell'apprendimento.

Come ragionava il progettista

Analisi del processo mentale che ha portato a una determinata soluzione tecnica.

L'obiettivo è insegnare il metodo, non soltanto il codice.

Dietro le quinte

Approfondimenti sulle tecniche utilizzate nei grandi videogiochi dell'epoca.

Lo scopo non è recensire i giochi, ma comprenderne le soluzioni progettuali.

Curiosità

Aneddoti storici, riferimenti culturali, citazioni di riviste, manuali e protagonisti dell'epoca.

\begin{tikzpicture}
\path  (0,0.054) --  (17.001,0.054) --  (17.001,0) --  (0,0) --  (0,0.054) -- cycle;
\end{tikzpicture}


\section{Materiale iconografico}
Le immagini non devono limitarsi a illustrare i concetti.

Devono raccontare un'epoca.

Lo stile grafico dovrà richiamare il celebre manuale del Commodore 64 senza copiarlo.

Le illustrazioni saranno reinterpretazioni moderne dell'immaginario Commodore.

La mongolfiera diventerà un drone.

Gli schemi saranno realizzati con un'estetica coerente con quella degli anni Ottanta.

Le fotografie dovranno alternarsi a diagrammi, schermate, sprite, mappe di memoria e figure esplicative per evitare la
monotonia tipica dei manuali tecnici.

\begin{tikzpicture}
\path  (0,0.054) --  (17.001,0.054) --  (17.001,0) --  (0,0) --  (0,0.054) -- cycle;
\end{tikzpicture}


\section{Materiale allegato}
Il libro sarà accompagnato dal progetto completo del videogioco.

Ogni capitolo corrisponderà a una versione realmente funzionante del codice.

L'idea è quella di permettere al lettore di ripercorrere tutta l'evoluzione del progetto.

Qualora il progetto editoriale lo consenta, il materiale potrà essere distribuito anche su un supporto fisico che
richiami la storica musicassetta del Commodore, ad esempio una scheda SD racchiusa in un contenitore a forma di
cassetta.

\begin{tikzpicture}
\path  (0,0.054) --  (17.001,0.054) --  (17.001,0) --  (0,0) --  (0,0.054) -- cycle;
\end{tikzpicture}


\section{Regole stilistiche}
La scrittura deve essere quella di un vero libro.

Mai quella di una conversazione.

Mai quella di un manuale scolastico.

Mai una sequenza di definizioni.

Prediligere periodi ampi, narrativi e scorrevoli.

Gli elenchi devono essere ridotti allo stretto necessario.

Alternare continuamente racconto, spiegazione, immagini, laboratori, esperimenti e approfondimenti per mantenere un
ritmo di lettura elevato.

Ogni capitolo deve avere una propria identità narrativa.

\begin{tikzpicture}
\path  (0,0.054) --  (17.001,0.054) --  (17.001,0) --  (0,0) --  (0,0.054) -- cycle;
\end{tikzpicture}


\section{La promessa al lettore}
Al termine della lettura il lettore non dovrà soltanto aver imparato a programmare il Commodore 64.

Avrà costruito un videogioco completo.

Avrà imparato il BASIC e l'Assembly.

Avrà compreso come ragionavano i programmatori degli anni Ottanta.

Avrà scoperto il funzionamento interno del VIC-II, del SID e dell'intera architettura del Commodore 64.

Ma soprattutto avrà ritrovato quella sensazione che molti ricordano ancora oggi: il momento in cui, davanti a una
schermata blu con un cursore lampeggiante, tutto sembrava possibile.


\bigskip
\end{document}
